\chapter{Acumulación mínima}

\section{Registro acumulativo}
\begin{definition}
Existe acumulación mínima cuando una sucesión de diferencias legibles puede conservarse como registro estructural sin colapsar de inmediato en indistinción.
\end{definition}

\begin{observation}
La acumulación mínima no implica todavía memoria persistente ni estabilidad larga. Solo fija la posibilidad de que lo leído no desaparezca instantáneamente como si no hubiera sido registrado.
\end{observation}

\section{Acumulación sobre trayectoria}
\begin{definition}
Una trayectoria mínima admite acumulación si sus pasos estructurales pueden incorporarse sucesivamente a un registro compatible con la resolución \eps\ y con la ventana \Win.
\end{definition}

\begin{observation}
Esta cláusula no exige todavía una ley cuantitativa completa de acumulación. Solo declara que el tránsito puede dejar rastro legible dentro del dominio local.
\end{observation}

\section{Acumulación y coherencia}
\begin{definition}
La acumulación mínima es compatible con la coherencia local cuando el registro de diferencias no destruye la legibilidad estructural de la configuración o trayectoria acumulada.
\end{definition}

\begin{observation}
No toda acumulación preserva coherencia. Esta definición solo admite acumulación cuando el registro sigue siendo tratable sin contradicción inmediata dentro del dominio declarado.
\end{observation}

\section{Acumulación y oposición}
\begin{definition}
Una acumulación es compatible con la involución \Iop\ si el registro acumulado puede ser invertido sin perder de manera inmediata su gramática estructural básica.
\end{definition}

\begin{observation}
Esta compatibilidad no fuerza simetría perfecta entre acumulación e inversión. Solo exige que la oposición siga siendo legible como transformación del registro.
\end{observation}

\section{Condición mínima de persistencia}
\begin{definition}
Se dice que una diferencia posee persistencia mínima si permanece registrable al menos a través de un paso adicional de lectura dentro de la resolución declarada.
\end{definition}

\begin{observation}
La persistencia mínima es una condición débil. Su función aquí es distinguir entre diferencia instantánea no retenida y diferencia que alcanza un primer umbral de conservación.
\end{observation}

\section{Cláusula de alcance}
\begin{observation}
La acumulación mínima fija una posibilidad elemental de conservación estructural. No valida por sí sola memoria fuerte, invariancia de largo alcance ni régimen dinámico estabilizado.
\end{observation}

